\section{آنچه داده گفت}

\begin{frame}{عدم‌دریافت یک تصمیم غذایی نیست، غیبت از دانشگاه است}
  \begin{itemize}
    \item اگر کسی ناهارش را نگیرد، احتمال نگرفتن شامش $3.2$ برابر می‌شود
    \item فرضیه‌ی «ترجیح غذایی شخصی» رد شد: پایداری جفت (فرد، غذا) $0.24$ در برابر پایداری
      خود فرد $0.43$
    \item \textbf{پیامد:} تغییر منو ابزار کاهش هدررفت نیست
  \end{itemize}
  \begin{figure}
    \centering
    \includegraphics[width=0.55\textwidth]{slides_08_dinner_given_lunch}
  \end{figure}
  \note{شاهد: \lr{F57}، \lr{F55}، \lr{F66}.}
\end{frame}

\begin{frame}{$83\%$ نوسان از یک شوک مشترک روزانه می‌آید}
  \begin{columns}
    \begin{column}{0.5\textwidth}
      \begin{itemize}
        \item تجزیه‌ی واریانس نرخ سلول: $7\%$ نویز نمونه‌گیری، $10\%$ ترکیب رزروکنندگان،
          $83\%$ شوک مشترک روز
        \item در $99.5\%$ جفت‌سلف‌ها این شوک هم‌جهت است — سراسری دانشگاه است نه محلی
        \item خبر خوب: با اطلاعات مجاز در لحظه‌ی برش تا $R^2=0.60$ پیش‌بینی‌پذیر است
      \end{itemize}
    \end{column}
    \begin{column}{0.5\textwidth}
      \begin{figure}
        \centering
        \includegraphics[width=\textwidth]{slides_09_variance_donut}
      \end{figure}
    \end{column}
  \end{columns}
  \begin{figure}
    \centering
    \includegraphics[width=0.62\textwidth]{report_14_day_shock_fa}
  \end{figure}
  \note{پیامد: مسئله در سطح «روز دانشگاه» مدیریت می‌شود، نه سطح فرد. شاهد: \lr{F59}، \lr{F60}، \lr{F61}.}
\end{frame}

\begin{frame}{روزهایی که بد به‌نظر می‌رسند، روزهایی نیستند که پول می‌سوزانند}
  \begin{itemize}
    \item همبستگی حجم رزرو با نرخ منفی است ($-0.35$) ولی با هدررفت مطلق شدیداً مثبت
      ($+0.75$)
    \item فهرست «بدترین روزها» با معیار نرخ و با معیار پرس فقط ۴ عضو از ۱۵ عضو مشترک دارند
  \end{itemize}
  \begin{figure}
    \centering
    \includegraphics[width=0.62\textwidth]{R5_q43_volume_vs_rate}
  \end{figure}
  \note{مهم‌ترین حرف مدیریتی ارائه. شاخص پایش سلف باید «پرس» باشد نه «درصد»، وگرنه مدیریت
  دقیقاً روزهای اشتباه را رصد می‌کند. شاهد: \lr{F80}، \lr{F81}، \lr{F54}.}
\end{frame}

\begin{frame}{هدف‌گیری دانشجویان پرغیبت جواب نمی‌دهد}
  \begin{itemize}
    \item سقف نظری هر سیاستی که $10\%$ افراد را هدف بگیرد، $38\%$ هدررفت است
    \item هدف‌گیری بر پایه‌ی مشخصات جمعیتی به $8.7\%$ می‌رسد — بدتر از انتخاب تصادفی
      ($10\%$)، چون افراد پرنرخ عمدتاً کم‌رزرواند
  \end{itemize}
  \begin{figure}
    \centering
    \includegraphics[width=0.62\textwidth]{report_12b_lorenz_targeting_fa}
  \end{figure}
  \note{پیامد: هر سیاست جریمه یا محدودیت بر پایه‌ی مقطع، جنسیت یا دوره مبنای داده‌ای ندارد.
  این اسلاید یک ایده‌ی بد را پیش از تولدش می‌کشد. شاهد: \lr{F86}، \lr{F74}، \lr{F73}.}
\end{frame}

\begin{frame}{تقویم سیگنال رایگان می‌دهد}
  \begin{itemize}
    \item روز قبل از تعطیلی: $3.1$ واحد درصد بالاتر، حدود $39\%$ نسبی
    \item چهار روز اول بازگشت پس از وقفه‌ی چهارروزه یا بیشتر: $22\%$ نسبی بالاتر — بدترین
      روز، روز \textbf{دوم} است نه اول
    \item شام رمضان تقریباً دو برابر
  \end{itemize}
  \begin{columns}
    \begin{column}{0.5\textwidth}
      \includegraphics[width=\textwidth]{report_07_pre_holiday_fa}
    \end{column}
    \begin{column}{0.5\textwidth}
      \includegraphics[width=\textwidth]{r4_2_return_curve_fa}
    \end{column}
  \end{columns}
  \note{همه‌ی این‌ها ماه‌ها قبل معلوم‌اند و اعمالشان هزینه‌ی صفر دارد، حتی بدون هیچ مدلی.
  شاهد: \lr{F19}، \lr{F70}، \lr{F22}.}
\end{frame}

\begin{frame}{جغرافیا}
  \begin{itemize}
    \item تهران $8.7\%$ در برابر رضوانشهر $3.1\%$، با تفکیک تقریباً کامل بین پردیس‌ها
    \item سازوکارش رفتار افراد است نه ترکیب سلف‌ها — در پردیس دورافتاده سلف تنها گزینه‌ی
      غذایی است
    \item \textbf{پیامد:} تمرکز منابع روی تهران
  \end{itemize}
  \begin{figure}
    \centering
    \includegraphics[width=0.62\textwidth]{report_02_city_effect_fa}
  \end{figure}
  \note{اگر وقت اضافه داشتی یا استاد پرسید، ماجرای آلودگی هوا را بیاور (\lr{F25}، اسلاید پشتیبان
  پ۱۱): اثر خام قوی بود ($r=+0.28$) و پس از کنترل سلف و وعده کاملاً صفر شد، چون تهران هم‌زمان
  بالاترین آلودگی و بالاترین نرخ را دارد. شاهد: \lr{F12}، \lr{F14}.}
\end{frame}
